\section{Closed four-dimensional toric-code estimates}\label{app:code-estimates}
Theorem~\ref{thm:exponential} uses a closed four-dimensional toric code on each fibre. We give the sweep predecessor rule, the grouped contour argument and the deterministic recovery criteria here. Lemma~\ref{x:exponential} reuses these arguments on a closed side-\(L\) code. Appendix~\ref{sec:folded} gives the folded code; Table~\ref{q:links} in Section~\ref{q:link-section} gives the common link table used by Theorem~\ref{thm:main}. Theorem~\ref{thm:main} also uses the grouped-contour estimate of Lemma~\ref{q:frontier} in Section~\ref{sec:controller} and Lemma~\ref{d:sources}.

\subsection{Sweep predecessors}

We use a variant of the sweep rule of Kubica and Preskill
\cite{KubicaPreskill2019}; for other local toric-code decoders see \cite{BDMT2017}. Simultaneously at every vertex, read the
positive incident edges of the syndrome present before the sweep; if
their number is even, pair them in a fixed order and flip the
corresponding positive plaquettes, and if it is odd, do nothing. This
is the positive sweep; reflection reverses all
directions. Fixed Clifford extraction and local Boolean record
processing implement the rule. Run \(400\) reflected sweeps per CSS
sector on the code complex (the virtual complex for the folded circuit).

In each round, for the bit sector followed by the phase sector, one
Clifford extraction writes the syndrome into fresh records, the snapshot
\(\sigma_0\). Local Boolean processing on these records alone computes
400 reflected sweeps by
\(\sigma_{j+1}=\sigma_j+\partial P_j\) for \(0\le j<400\), where
\(P_j\) is the plaquette set that sweep \(j+1\) selects from
\(\sigma_j\), and stores \(P=P_0+\cdots+P_{399}\). Only then are the
plaquette flips of \(P\) applied. The syndrome present before sweep
\(j+1\) means \(\sigma_j\); no sweep reuses an earlier sweep's decisions.
The record processing has range at most 400 code cells.
Track both sector chains at the common boundaries of full rounds,
just before the bit-sector extraction. The catalogue \(\Lambda_t\)
includes the allowed insertions and erroneous uses throughout the full
round, as in Lemma~\ref{q:round-shadow}.

\begin{lemma}[Sweep predecessors]\label{q:sweep}
For one positive sweep, every vertex of the new syndrome has, separately
for each of the four functions \(x_j\), an old syndrome predecessor with
nonnegative gain in that function. It also has an old predecessor with
strictly positive gain in \(-\sum_jx_j\). The reflected sweep has the
corresponding reversed-coordinate property. After \(M\) reflected
sweeps the gains for the walls \(-x_j\) and \(\sum_jx_j\) are at least
\(0\) and \(M\), respectively.
\end{lemma}
\par\noindent\emph{Proof idea.} Trace a new syndrome corner to the plaquette correction that created it. The old incident edges give separate wall predecessors, and the correction root gives the strict sum gain.

\begin{proof}
Suppose a new syndrome vertex \(v\) was absent from the old syndrome.
It is a corner of a corrected plaquette based at an old syndrome vertex
\(u\). If its \(j\)-coordinate was raised, the old positive edge in
direction \(j\) used at \(u\) has the raised endpoint as an old syndrome
vertex. That endpoint supplies the \(x_j\) predecessor. If the coordinate
was not raised, use \(u\). The correction root \(u\) has smaller
coordinate sum than \(v\), and supplies the last predecessor.

An old vertex that remains a syndrome vertex supplies its own first
four predecessors. If it has an old negative incident edge, that
edge's lower endpoint supplies the last one. Otherwise every old
incident edge is positive. Their number is even because the syndrome
is a cycle. The vertex's own correction then cancels all of them.
Any incident edge remaining after the parallel correction must be
toggled by a correction based at a vertex of smaller coordinate sum;
that correction root is an old syndrome vertex. This proves the strict
last gain. Cancellations between parallel corrections do not change
this argument. Reflection and iteration give the stated bounds. The
dual lattice differs by a fixed half-cell displacement and has the
same proof.
\end{proof}

\begin{lemma}[Grouped derivations]\label{q:frontier}
Suppose a finite acyclic backward derivation has \(k\) balanced
potentials. For each potential every positive support has a selected
parent with gain at least \(d_0>0\), and every selected edge loses at
most \(K_0\). Allow source-free zero-gain corridors, followed by a
positive support, and bounded entry and terminal links, at most one
of each per arm. Group terminal appearances by source label. For group
\(i\), let \(\Omega_i\) bound the total potential oscillation and let
\(K_{{\rm term},i}\) bound the loss of a terminal link. Every root has a
connected witness with \(g\) distinct true source labels. If \(E\) counts
positive edges with arm multiplicity and \(K_{\rm ent}\) bounds the entry
loss on one arm, then
\[
 d_0 E\le k(d_0+K_0)(g-1)+kK_{\rm ent}g+
 \sum_{i=1}^g(\Omega_i+kK_{{\rm term},i}).
\]
With uniform terminal loss and oscillation bounds, the projected witness
has at most \(C_{\rm w}g\) edges, where \(C_{\rm w}\) is computable from
\(k,d_0,K_0,K_{\rm ent}\), the fixed typed stencils and these bounds.
\end{lemma}
\par\noindent\emph{Proof idea.} Keep one arm endpoint of each charge in every overlap component of the frontier. A split requires new arms, and balancing their potential increments bounds the total length by the terminal losses and oscillations.

\begin{proof}
The argument refines Toom's contour method
\cite{BermanSimon1988,Gacs2001,Toom1980} to grouped terminal labels.
For a derivation vertex \(v\), let \(T(v)\) be the set of terminal
labels reachable from it. Begin with the frontier consisting of the
root and expose a frontier nonterminal by replacing it with all its
selected parents. Partition the frontier by connected components of
the intersection graph of its \(T\)-sets. New sets are subsets of the
exposed set, so exposure cannot merge previous components.

If one component splits into \(j\) daughters, each daughter contains
a selected parent of the exposed vertex. Indeed, any remaining old
frontier set previously attached to the exposed set intersects a
terminal set of at least one selected parent; any surviving connection
not using the exposed set remains. For each charge, keep one arm
endpoint in every current component. Initially follow a charge-good
arm from the root to a terminal. At a split the existing endpoint
belongs to one daughter. Add \(j-1\) source occurrences at the exposed
vertex. For that charge route their first edges into the daughters
missing an endpoint, then continue by charge-good edges. The chosen
parent's whole reachable terminal set is in that daughter, so the
continuation cannot leave it. The assignments for different charges
may differ.

At the terminal frontier each distinct terminal label receives
exactly one arm endpoint per charge. Join the starts at each source
occurrence and the endpoints at each terminal label, and retain the
undirected arm component containing the original root. In this
component every source starts one arm of each charge and every
terminal receives one of each charge. Therefore \(g\) retained
terminal labels correspond to \(g\) retained source occurrences.
Exactly \(g-1\) occurrences were added, so at most \(k(g-1)\) first
edges fail to be charge-good.

Sum potential increments along all arms. Balanced potentials cancel
at each source. At terminal group \(i\), cancellation costs at most
\(\Omega_i\). There are \(kg\) arms, with at most one entry link per arm
and one terminal link of each charge at group \(i\). Summing these
losses gives the displayed bound. Under uniform loss and oscillation
bounds, \(E=O(g)\). A zero corridor cannot split an overlap component. After
projection each such corridor contributes a fixed number of links,
charged to its following positive support; root and terminal
exceptions cost \(O(g)\). The entry and terminal links also cost
\(O(g)\). Projection preserves connectedness and can only identify
edges, giving the asserted witness.
\end{proof}

\subsection{One recovery for every permitted Pauli assignment}

We first record the deterministic topological criterion. Use the shortest period of the intended closed code torus.

\begin{lemma}[Local filling invariant]\label{q:homology}
Let \(\ell_{\min}\) be the shortest period of the code torus and set \(r=\ell_{\min}/20\).
Starting from a code state, if every syndrome component in both CSS
sectors has coordinate diameter less than \(r\) at every serialized
elementary plaquette Pauli update, a fixed final filling decoder
recovers the entire logical state and its reference exactly.
\end{lemma}
\par\noindent\emph{Proof idea.} Fill each small lifted syndrome component in a contractible box. A single plaquette update changes the corrected chain by a boundary, preserving its homology class.
\begin{proof}
A syndrome component is itself a \(1\)-cycle. Its small lift lies
in a contractible box, so fix a \(2\)-chain filling there by finite
Gaussian elimination and a deterministic box/lift convention. Set
\[
 h(e)=\left[e+\sum_\nu C_\nu(\partial e)\right]
 \in H_2(\mathbb T^4;\mathbb F_2).
\]
Under one plaquette update only components meeting that plaquette
can change. The affected old and new components and their fillings
lie in a common box of width at most \(8r+O(1)<\ell_{\min}\), after
enlarging the fixed size threshold. Unaffected fillings cancel.
The difference of the two corrected chains is closed in that box,
hence a \(3\)-boundary. Thus \(h(e)\) is unchanged. It starts at
zero. Applying the same argument on the dual cellulation shows that
the final Pauli acts by stabilizers in both sectors, including on
reference-entangled states.
\end{proof}

\begingroup
\renewcommand{\thetheorem}{\thesection.3a}
\renewcommand{\theHtheorem}{\thesection.3a}
\begin{lemma}[Round shadowing]\label{q:round-shadow}
Track both sector chains at common full-round boundaries, just before
the bit-sector extraction. Fix a round \(t\) and a sector, and let
\(e_t,e_{t+1}\) be its Pauli-frame 2-chains at those boundaries.
Let \(\Lambda_t\) include the catalogue locations of allowed insertions
and erroneous uses throughout that full round: inserted Paulis,
erroneous records, and erroneous stored decisions or mask uses.
Let \(R_*\) bound the code-lattice radius of each propagated local
contribution through every remaining stage of that full round, namely
to this sector's extraction, through its record processing, and to the
next common boundary, uniformly over both sectors.
Let \(K_*\) bound the number of sheet patches per location, let
\(d_\square\) be the diameter of a plaquette, and put
\[
 R_{\rm blk}=400d_\square+2R_*.
\]
For every serialization of the round's Pauli updates and every
intermediate chain \(e\), every vertex of \(\partial e\), and every
vertex of \(\partial e_{t+1}\), lies within distance \(R_{\rm blk}\)
of \(\partial e_t\cup\Lambda_t\). For the folded circuit the same
holds in each sheet's virtual chart, with the \(K_*\) sheet patches
of each location.
\end{lemma}
\addtocounter{theorem}{-1}
\endgroup
\begin{proof}
At this sector's extraction write \(\sigma_0=\partial e_t+\rho\),
where \(\rho\) includes erroneous records and data Paulis from the
preceding stages of the same round. Each such propagated local
contribution lies within \(R_*\) of its location in \(\Lambda_t\).
Ideal operations of the other CSS sector leave this sector's syndrome
unchanged. An erroneous intermediate sweep record or data Pauli
enters only within \(R_*\) of its location in \(\Lambda_t\).
A vertex of
\(\sigma_{j+1}\) absent from \(\sigma_j\) is a corner of a plaquette
of \(P_j\). That plaquette is selected at a vertex of \(\sigma_j\)
by the proof of Lemma~\ref{q:sweep}, or, if its stored decision is
erroneous, lies within \(R_*\) of \(\Lambda_t\). Induction therefore
places every vertex of \(\sigma_j\) within \(jd_\square+R_*\) of
\(\sigma_0\cup\Lambda_t\). An intermediate chain is
\(e=e_t+F'+P'\), where \(F'\) consists of inserted Paulis of this
round and \(P'\) consists of applied plaquettes of \(P\). A vertex
of \(\partial e\) lies in \(\partial e_t\), in an inserted Pauli's
boundary within \(R_*\) of \(\Lambda_t\), or on a plaquette of
some \(P_j\), within \(d_\square\) of \(\sigma_j\), hence within
\((j+1)d_\square+2R_*\) of \(\partial e_t\cup\Lambda_t\). Since
\(j\le399\), these distances are at most \(R_{\rm blk}\).
The final chain is included. The argument is local in each virtual
chart, so it applies separately to every sheet patch.
\end{proof}

A fixed exogenous classical history determines a catalogue of allowed
Pauli-insertion supports: physical-noise supports, erroneous record
or mask uses, controller erasures, and each affected sheet image of
a wrong program gate. The nonunitary instruments of the intended
processor have one common unconditional Stinespring dilation. Only
unitaries differ between actual and intended programmed circuits.
For the folded circuit choose \(r_{\rm d}\) equal to a sufficiently small fixed fraction of \(m\) from Appendix~\ref{sec:folded}. An admissible dangerous certificate is a
connected typed-source certificate from the proof of Lemma~\ref{q:spatial}
spanning at least a fixed multiple of \(r_{\rm d}\), with allowed
predecessor, program-use and terminal annotations. Define
\(\mathcal G\) to mean that none exists. This definition precedes
and is independent of any choice of Pauli values at the catalogue.

The folded circuit uses the source derivations of Appendix~\ref{sec:folded}.
For the six-dimensional circuit, \(\mathcal G\) is the complement of the
exceptional event in Lemma~\ref{v:bridge}. Its proof supplies the corresponding
derivations, admissible certificates and probability bound.

\begin{lemma}[Universal certificate implication]\label{q:universal}
For a fixed exogenous classical history, every Pauli assignment at
the allowed catalogue locations that changes corrected homology in
either CSS sector yields an admissible dangerous certificate.
Consequently, on \(\mathcal G\), the same intended-complex filling
decoder corrects every insertion history on the entire logical
codespace, including its reference.
\end{lemma}
\par\noindent\emph{Proof idea.} Fix an arbitrary Pauli history and use stabilizer extraction to track its syndrome without reference to the logical state. A first change of corrected homology forces a large component, whose predecessor graph yields the certificate.

\begin{proof}
Fix an arbitrary complete insertion assignment \(h\). In the ideal
inserted network the data are always a Pauli frame on a codeword.
Indeed, Clifford extraction conjugates Pauli observables to Pauli
observables, and conjugation by an inserted Pauli changes only their
signs. A stabilizer measurement therefore has a deterministic record
for this frame, independent of the stored logical state. Record
Paulis flip fixed basis records; the local Boolean matching then
computes deterministic correction choices. Its work remains basis
classical. In the common instrument dilation the associated
measurement environment vector is logical-independent. This proves
the assertion for the entire successive history by induction.

In a fixed extraction/sweep block, a syndrome vertex has the ideal
\(400\)-sweep predecessors of Lemma~\ref{q:sweep} unless its ideal
backward cone contains an inserted or erroneous location. In that
case terminate its quantum arm at the last such location in the
block. The remaining ideal propagation has a fixed radius in that
sheet. Attach the corresponding per-sheet link, then its true
profile-use or selector derivation, or its fresh hardware source.
Record errors have the same bounded local correction influence.
The derivation is finite and its classical entries never return to
quantum types. In the folded case Table~\ref{q:links} supplies all increments,
including repeated predictor blocks. In six dimensions the typed source label catalogue
in the proof of Lemma~\ref{v:bridge} supplies the links and the grouped derivations.

A Pauli term may have several separated sheet patches. Let \(K_*\) and
\(R_*\) be as in Lemma~\ref{q:round-shadow}; they bound the number and
the radii of these patches.
Choose \(r_{\rm d}\) small enough that any union of \(K_*\) patches
linked at distances \(2(r_{\rm d}+R_*)\) fits in a contractible box much
narrower than \(\ell_{\min}\). Combine such nearby patches into
one box. Separated boxes can be applied successively: if the
before and after syndromes have only components smaller than
\(r_{\rm d}\), no such component joins two separated boxes. An
intermediate connected piece therefore belongs to the before
syndrome near an unprocessed box or to the after syndrome near a
processed box, together with the bounded patch neighbourhoods.
Inside a combined box serialize its plaquettes. Every intermediate
component is contained in a fixed multiple of that contractible
box. Taking the fixed fraction defining \(r_{\rm d}\) smaller if needed
puts these diameters below the barrier in
Lemma~\ref{q:homology}.

It follows that a change of corrected homology forces a
macroscopic syndrome component either at a tracked block boundary
or in such a local serialization. Map each vertex of the first macroscopic intermediate component \(K\)
to a shadow root in \(\partial e_t\cup\Lambda_t\) within
\(R_{\rm blk}\), by Lemma~\ref{q:round-shadow}. This typed link has
fixed range, decreases physical time by at most one round, and loses
at most a fixed amount of each potential; it is the intermediate-stage
spur. A tracked root carries its grouped witness from
Lemma~\ref{q:frontier}. Its 400-sweep link (Table~\ref{q:links}) is
Lemma~\ref{q:sweep} applied to its own round's exact recurrence
\(\sigma_0,\ldots,\sigma_{400}\), terminated at the last location
of \(\Lambda\) in its backward cone. A root in \(\Lambda_t\)
is a single-source witness. Adjacent vertices of \(K\) have shadow
roots at distance at most \(2R_{\rm blk}+1\), so the shadow roots
form a chain spanning at least \(\operatorname{diam}K-2R_{\rm blk}\).
In the ball covering use \((4+2R_{\rm blk}+1)\)-fold enlargements.
The constants \(C_{\rm raw},C_{\rm all},\Delta_*,\Delta_{\rm all},
b_0,b_{\rm span}\) change only through \(R_{\rm blk}\), fixed
with the catalogue. Every local patch is
anchored at its allowed typed insertion source; no long virtual
edge is inserted between distinct sheet images. Apply the
predecessor construction just proved and Lemma~\ref{q:frontier},
then the ball and connector construction in the proof of Lemma~\ref{q:spatial}
for the folded circuit or in the proof of Lemma~\ref{v:bridge} for the
six-dimensional circuit, to obtain the dangerous certificate.

All these constructions use only positions, finite roles, and true
terminal events. The argument applies to the arbitrary assignment
\(h\) with which we began. The certificate count therefore controls
the union over every assignment. On \(\mathcal G\) no assignment
can change homology, and the one fixed filling decoder returns every
logical state exactly by Lemma~\ref{q:homology}.
\end{proof}

\begin{lemma}[Common recovery and coherent span]\label{q:span}
Let \(V\) encode a logical system \(A\), also writing \(V\) for its encoding channel,
and suppose one recovery
\(\mathcal D\) satisfies \(\mathcal D\Phi_hV=\operatorname{id}_A\)
for every allowed Pauli-inserted history in a common unconditional
instrument network. Then it also corrects coherent unitary changes
whose Pauli expansions use those locations. If changed gates couple
to an ideal-decoupled outside system \(E\), the corrected dilation
has the form \(I_A\otimes K_E\), with \(K_E\) an isometry. Hence
logical/reference marginals are preserved even for initially
correlated logical and outside inputs.
\end{lemma}
\par\noindent\emph{Proof idea.} Exact recovery of each inserted history factors the recovered logical output from its environment in a common dilation. Expand the actual gates in Pauli operators and sum the factors, including the outside-system operators.

\begin{proof}
Choose the common history dilations \(\Omega_h\) and one recovery dilation
\(R\). Exact identity recovery on a logical basis gives
\[
 R\Omega_hV|a\rangle=|a\rangle\otimes|s_{h,a}\rangle,
\]
since the recovered marginal is pure. Preservation of the
off-diagonals forces
\(\langle s_{h,b}|s_{h,a}\rangle=1\), so these unit vectors are
all equal, including their phase. Therefore
\(R\Omega_hV=I_A\otimes|s_h\rangle\).

For each changed unitary write
\(U_{\rm act}U_{\rm ideal}^*=\sum_P\alpha_P P\).
Expanding the finite interleaved product in the common dilation gives
\(\Omega_{\rm act}V=\sum_h\alpha_h\Omega_hV\). Applying \(R\) yields
\[
 R\Omega_{\rm act}V=I_A\otimes\sum_h\alpha_h|s_h\rangle.
\]
The actual dilation is isometric, so the remaining vector has norm
one. The full coherent sum, including all cross terms, has exact
identity recovery.

For an outside system instead expand each difference as
\(\sum_P P_{\rm active}\otimes B_P\). The intended active/outside
factorization allows all outside operations and their dilation
baths to remain in these operator-valued coefficients. The same
calculation gives
\[
 (R\otimes I)\Omega_{\rm act}(V\otimes I_E)=I_A\otimes K_E.
\]
Isometry implies \(K_E^*K_E=I_E\). Its partial trace is a channel
on the environment, so the logical/reference marginal is unchanged
for every joint input on \(A,E\), and a reference. Active check
ancillas and records remain in \(\Omega_h\); only ideal-decoupled
outside factors enter \(E\).
\end{proof}

\subsection{Backward increments for common witness links}\label{q:link-section}
For a local lifted virtual chart set \(d=(58,58,58,57)\) and
\begin{equation}\label{q:walls}
 F_j=-(q_j+d_jt)\ (1\le j\le4),\qquad
 F_5=\sum_{j=1}^4(q_j+d_jt).
\end{equation}
They sum to zero. Table~\ref{q:links} lists backward increments
\(F(\text{parent})-F(\text{root})\). It includes links created by
classical gathering and predictor recentring. Write
\(\delta_i\in\{-1,0\}\) and \(T_{\rm hist}=1024\). Here \(u_j\) is the \(j\)-th component of the selected replica displacement in the selector or profile box.

\begin{table}[!htbp]
\centering\small
\renewcommand{\arraystretch}{1.1}
\caption{Backward increments \(F(\text{parent})-F(\text{root})\) of the five balanced potentials \(F_1,\ldots,F_5\) in \eqref{q:walls}. The witness graphs in Lemmas~\ref{d:sources}, \ref{q:universal} and~\ref{q:spatial}, and Proposition~\ref{c:oriented-source}, use these links. The five entries are in potential order. The good entries for different charges may use different parents.}
\label{q:links}
\begin{tabular}{>{\raggedright\arraybackslash}p{.25\textwidth}>{\raggedright\arraybackslash}p{.68\textwidth}}
\toprule
Link & Five increments, or their justified lower bounds\\
\midrule
Selector support &
\((58-u_1,58-u_2,58-u_1,57-u_2-\delta_i)\),\newline
\(\Delta F_5=2u_1+2u_2+\delta_i-231\).
Charge-specific good choices give \((43,43,43,42,88)\).\\\addlinespace[3pt]
Profile support &
\((58-u_1,58,58-u_3,57)\),\newline
\(\Delta F_5=u_1+u_3-231\).
\newline The \(320\)-separated boxes give \((43,58,43,57,89)\).\\\addlinespace[3pt]
\(400\) reflected sweeps &
\((58,58,58,57,169)\) as charge-specific lower bounds.\\\addlinespace[3pt]
Sampling-register copy corridor &
\((0,0,0,0,0)\) on its base characteristic.\\\addlinespace[3pt]
Gather then field consult &
\((0,0,0,+j,-j)\) for the gather \(i\mapsto i-j\);
\((0,0,0,-j,+j)\) for the consult back to the base.
Their compressed link has zero increment.\\\addlinespace[3pt]
Clean predictor entry &
\((0,0,0,T_{\rm hist}-1,-(T_{\rm hist}-1))\), followed by all \(T_{\rm hist}\) formal support levels.\\\addlinespace[3pt]
Same-site old/new star &
\(\pm(58,58,58,57,-231)\); total absolute oscillation \(462\).\\\addlinespace[3pt]
Quantum or program-use link &
For its fixed displacement \((\Delta q,\Delta t)\),
\((-\Delta q_j-d_j\Delta t)_{j=1}^4\) and
\(\sum_j\Delta q_j+231\Delta t\). Each sheet's allowed list is finite.\\
\bottomrule
\end{tabular}
\end{table}

\subsection{Finite-torus exponential storage}\label{q:closed-exponential}
\begin{lemma}[Finite-torus exponential storage]\label{x:exponential}
Consider the closed side-\(L\) four-dimensional plaquette toric code,
with the reflected-sweep processor of Lemma~\ref{q:sweep} in both
CSS sectors. An iid exogenous site-round source of sufficiently small
computable rate \(p_{\rm aux}>0\) may cause arbitrary faults in a
fixed finite influence catalogue. There are computable constants
\(A,f_0>0\) such that the probability of any dangerous
certificate through any finite time \(T\) is at most
\begin{equation}\label{x:exp-bound}
 A(T+A)L^4\cdot2^{-f_0L}.
\end{equation}
On its complement one fixed component-filling decoder recovers the
whole logical state and reference. Consequently this processor stores
a qubit with fidelity at least \(15/16\) through
\(2^{\lfloor c_{\exp}L\rfloor}\), for rational
\(0<c_{\exp}<f_0/4\) and all sufficiently large computable sizes.
\end{lemma}
\par\noindent\emph{Proof idea.} Use the sweep predecessors to attach a connected explanation to each syndrome vertex. A large component forces an explanation with many distinct sources, and excluding those explanations makes component filling valid for every allowed history.

\begin{proof}
There are six plaquette data qubits per site, with local extraction,
syndrome and correction work. Each sector uses the extraction and
400-step virtual recurrence specified above; the resulting Pauli mask is
stored before its gates execute. All
incidence layers and private controls are fixed. Begin with zero
syndrome and lift the finite-time dynamics periodically to \(\Z^4\).
Use the five balanced potentials
\[
 F_j=-(r_j+58t),\quad 1\le j\le4,\qquad
 F_5=\sum_j(r_j+58t).
\]
After 400 reflected sweeps, the backward negative potential gains are at
least 58 and the sum potential gain is at least \(400-4\cdot58=168\).
The intermediate-stage spurs of Lemma~\ref{q:round-shadow} have bounded losses, as constructed in the proof of Lemma~\ref{q:universal}. Group all
appearances of one lifted source under its single site-round name.
Lemma~\ref{q:frontier} gives a connected witness with at most
\(C_{\rm w}g\) edges for \(g\) distinct lifted source anchors.
This is a deterministic statement about a periodic field.

For a syndrome path of macroscopic span, enclose each root's complete
witness and its source anchors in a ball of radius \(O(g)\). Select
disjoint balls in decreasing radius. The
\((4+2R_{\rm blk}+1)\)-fold enlargements from
Lemma~\ref{q:round-shadow} and the proof of Lemma~\ref{q:universal}
cover the polygonal path, including edges between consecutive shadow roots. Choose a
simple connecting chain, join successive selected roots at the common
time, and join both original path endpoints to the first and last
selected roots. These links cost \(O(\sum g)\); a single large
selected ball still includes both endpoint links. Original disjoint
balls have disjoint lifted source sets. The combined graph has
\begin{equation}\label{x:raw-span}
 m\le C_{\rm raw}g,\qquad m\ge b_0L,
\end{equation}
where \(m\) is its edge count and \(b_0>0\) is the fixed span constant of this construction. All full space-time roles and canonical
source anchors are retained. Here \(R_{\rm edge}\) is a local edge-range
bound, chosen after \(R_{\rm blk}\) and including the already
constructed shadow links; it bounds each edge's spatial range.

Take \(q=\lfloor f_0L\rfloor\), reducing rational
\(f_0>0\) so that \(q\le m\) and \(R_{\rm edge}q<L/2\).
Double a spanning tree of the graph, obtaining a walk of length at
most \(2m\). Mark the first visit to each of its \(g\) source
anchors, and partition the walk into pieces of at most \(q\) edges.
There are at most \(2m/q+1\le3m/q\) pieces. One contains at least
\(gq/(3m)\ge q/(3C_{\rm raw})\) distinct marks. Assign a boundary
mark to one half-open piece, with the final endpoint in the last
piece. Each mark is assigned once and belongs to its piece.

This short connected walk has spatial diameter below \(L/2\).
Two distinct canonical lifted anchors in it cannot represent the
same physical site-round coin: a nonzero spatial deck translation
has a coordinate of magnitude at least \(L\), and time is not
quotiented. Thus its marks denote distinct physical iid sources.
No independence is assigned to the full old lifted graph.

If \(\Delta_*\) bounds the enlarged role graph's degree, including
a stay-put option, count short walks and their source subsets.
There are at most \(A_0(T+A_0)L^4\) physical starts and
\(2^{q+1}(\Delta_*+1)^q\) annotated walks. Given a start and its
edge sequence there is one lift. Discard annotations with repeated
physical source names. Independence is now valid and gives
\[
 \Pr(\text{dangerous certificate by }T)
 \le2A_0(T+A_0)L^4
       [2(\Delta_*+1)p_{\rm aux}^{1/(3C_{\rm raw})}]^q.
\]
For example choose the rational
\(p_{\rm aux}=[1/(8(\Delta_*+1))]^{3C_{\rm raw}}\).
With the actual source rate at most \(p_{\rm aux}\), the bracket is at most \(1/4\). Increasing the prefactor for the
floor in \(q\) gives \eqref{x:exp-bound}.

It remains to connect certificates to logical failure. Serialize
bounded local Pauli updates of each actual layer, keeping the stored
mask from the exact 400-step recurrence fixed. Use the shadow-root map
of Lemma~\ref{q:round-shadow} as in the proof of Lemma~\ref{q:universal}. Lemma~\ref{q:homology}, now with shortest side
\(L\), says that corrected homology cannot change while all affected
components have small contractible lifts. At the first barrier
crossing a bounded update merges only a fixed number of small
components. Choose the barrier fraction small enough that this first
component still has an injective lift. Its extreme vertices supply
the path used in \eqref{x:raw-span}. The predecessor history can be
arbitrarily old, but is finite at that crossing time. Intermediate
stages are included, so a closed error surface with zero endpoint
syndrome cannot bypass this argument.

Define the good event by absence of all such geometric certificates,
before assigning Pauli values. The predecessor and first-crossing
proof applies to every assignment in the fixed catalogue. The same
filling decoder therefore corrects their whole union. For general
local channels use Lemma~\ref{q:span}'s operator-valued expansion in
the common unconditional instrument dilation. It preserves the full
logical/reference state. Initialize one of the six logical pairs with
the unknown qubit and fix the other five. The same encoder works
throughout the horizon. Substituting
\(T=2^{\lfloor c_{\exp}L\rfloor}\) in \eqref{x:exp-bound} gives
failure at most \(1/16\) above a computable threshold.
\end{proof}

\subsection{Origin of recurring constants}\label{q:constants-origin}
\begin{itemize}
\item \(768=3\cdot16^2\) counts the voting replicas in the three boxes of \eqref{eq:c-boxes}. The threshold is \(512\). In the raw six-dimensional graph the nine Moore offsets give degree \(2\cdot9\cdot768=13\,824\), and the doubled walk count uses \(2\cdot13\,824^2=382\,205\,952\) (Lemma~\ref{v:density}).
\item The drift \(58\) in each protected coordinate balances the gains on the voting boxes. For each negative potential a replica outside its far box has \(s_j\le15\), giving \(3(58-15)=129\). A far-box replica has \(s_1+s_2\ge160\), giving \(3(160-2\cdot58)=132\) for the sum potential. The largest single-coordinate offset is \(175\), giving the worst loss \(3(175-58)=351\). An old and a new appearance differ by one step of this drift, giving total potential oscillation \(348+174+174=696\). The ordinary unshifted density proof instead uses gain \(130\) and loss \(350\) (Lemmas~\ref{c:coherence} and~\ref{v:density}, and Section~\ref{v:parameters}).
\item In the six-dimensional construction a complete query uses \(4096\) virtual levels. Distributing its entry loss of at most \(696\) over those levels leaves \(129-696/4096>128\) gain per edge (Section~\ref{v:parameters}).
\item A fork costs at most \(480\) per charge relative to a good edge. With three charges this gives \(1440\) per fork. In the density count, \(1440/130<12\) and \(1404\cdot12=16\,848\), with \(1/\delta=1404\); counting spurs and contours adds \(12\), giving \(16\,860\) (Lemma~\ref{v:density}).
\item The code uses \(400\) reflected sweeps. The drift costs \(4\cdot58\) on the sum potential, leaving \(400-4\cdot58=168\) in the six-dimensional proof (Lemmas~\ref{q:sweep} and~\ref{v:bridge}).
\end{itemize}
